\documentclass[11pt,oneside]{book}
\usepackage[T1]{fontenc}
\usepackage[utf8]{inputenc}
\usepackage[margin=1in]{geometry}
\usepackage{amsmath,amssymb}
\usepackage{microtype}
\usepackage[colorlinks=true,urlcolor=blue]{hyperref}
\newcommand{\courseweek}[1]{\textbf{#1}}
\setlength{\parindent}{0pt}
\setlength{\parskip}{6pt}
\title{Machine Learning in Engineering Science\\Sequential Data, LSTM, Attention, and Transformers}
\author{Chi-Hua Jonny Yu\\NCKU / LAiMM}
\date{Fall 2026 -- Week 7 -- 19 October 2026}
\begin{document}
\maketitle
\tableofcontents
\chapter{Sequential Data, LSTM, Attention, and Transformers}
\label{chap:sequences-attention}
\courseweek{Week 7: Monday 19 October 2026, 14:00--17:00.}

The three one-hour modules follow one question: how can a learned model use
earlier information to make a defensible prediction? Hour 1 develops recurrent
memory, Hour 2 develops query-dependent retrieval, and Hour 3 assembles and
tests a causal Transformer. About 90 minutes of explanation and worked
reasoning are combined with 90 minutes of prediction, guided computation,
experimentation, and explain-back across the three hours.

\section{Learning outcomes and evidence}
Students should be able to define a forecast origin and horizon; construct
windows without future leakage; trace every LSTM gate; compute attention
scores, weights, and a context vector; distinguish self-attention from
cross-attention; explain positional information, multi-head attention,
residuals, and feed-forward layers; and test a decoder's causal behavior.
The three independent CPU notebooks compare an LSTM against baselines,
verify attention arithmetic, and train and inspect a small Transformer.
All data are synthetic or self-contained. No paid API or GPU is required.
These are class exercises, not additional weekly graded assignments.

\section{Hour 1: ordered observations and recurrent memory}
\subsection{Order, sampling, and the prediction task}
A sequence is an ordered collection $x_1,\ldots,x_T$ in which order is part
of the problem definition. Sensor samples, process events, and language
tokens are examples. $(1,2,4,8)$ and $(8,1,4,2)$ have the same multiset and
mean but different transitions. A representation that discards order
cannot recover those transitions from the values alone. Equally spaced
samples and irregular events also require different sampling assumptions.
A sequence model does not automatically infer omitted time intervals.

In the example, the forecast origin is immediately before observing $x_t$.
With window length $L$,
\[X_t=(x_{t-L},\ldots,x_{t-1}),\qquad y_t=x_t.\]
The Python slice \texttt{signal[t-L:t]} excludes index $t$. Every input
index must be strictly less than the target index. This is a one-sample
horizon. A fixed-origin multi-step forecast cannot use later actual values
that would not yet have been observed.

With $B$ histories, $T$ time steps, and $F$ features at each step, batch-first
input has shape $(B,T,F)$. Many-to-one predicts one value per history;
many-to-many predicts outputs at multiple positions. Batch size is not
sequence length, and an ID is not necessarily a physical feature.

\subsection{Chronological splits and available information}
The synthetic sensor uses 1600 samples, window length 24, training targets
before index 960, validation targets 960--1279, and test targets 1280 onward.
This produces 936 training, 320 validation, and 320 test windows.
Target periods are disjoint. A test window may use earlier observed
validation or test sensor history. This is explicitly a rolling one-step
protocol, not a forecast of all future values from one initial origin.

Normalization learns statistics, so fit it on training observations only:
\[z_t=(x_t-\mu_{\mathrm{train}})/\sigma_{\mathrm{train}}.\]
Global means can leak future distribution information even without labels.
Random splits of heavily overlapping windows can produce nearly identical
training and test histories. For multiple devices or people, group separation
may be needed in addition to temporal separation.

\subsection{Baselines and error interpretation}
Persistence uses $\hat x_t=x_{t-1}$ and has no fitted parameters. Ridge
regression uses a regularized linear map of the flattened history window.
Select regularization on validation, then compare the same test targets:
\[
\operatorname{MAE}=\frac1N\sum_n|\hat y_n-y_n|,\quad
\operatorname{RMSE}=\sqrt{\frac1N\sum_n(\hat y_n-y_n)^2}.
\]
Report original target units and inspect residuals against time or operating
range. RMSE emphasizes large errors. Neither score alone reveals every
failure. A linear model can outperform an LSTM on a small, structured signal;
report the observed comparison rather than choosing the more complex name.

\subsection{RNN state and learning through time}
A basic recurrent network reuses one update at every time step:
\[h_t=\tanh(W_xx_t+W_hh_{t-1}+b).\]
Unrolling displays repeated applications of shared parameters. The state
is a learned summary, not necessarily an interpretable physical variable.
Backpropagation through time includes products of Jacobians:
\[\frac{\partial h_t}{\partial h_{t-k}}=J_tJ_{t-1}\cdots J_{t-k+1}.\]
Repeated contractions may shrink gradients, and repeated expansions may
amplify them. A scalar example $0.9^{100}$ illustrates one contraction
mechanism but not a complete nonlinear Jacobian product. Gradient clipping
limits large gradient magnitude; it does not itself solve vanishing gradients.

\subsection{The complete LSTM mechanism}
An LSTM separates cell memory $c_t$ from hidden output $h_t$. With
$u_t=[h_{t-1};x_t]$, a standard cell uses
\begin{align}
f_t&=\sigma(W_fu_t+b_f), & i_t&=\sigma(W_iu_t+b_i),\\
g_t&=\tanh(W_gu_t+b_g), & o_t&=\sigma(W_ou_t+b_o),\\
c_t&=f_t\odot c_{t-1}+i_t\odot g_t, &
h_t&=o_t\odot\tanh(c_t).
\end{align}
Here $\sigma(z)=1/(1+e^{-z})$ and $\odot$ is elementwise multiplication.
Each gate is a vector of coefficients, generally one per hidden dimension.
The forget gate retains old cell content; the input gate scales a signed
candidate write; the output gate exposes a filtered view of the updated cell.
The additive memory path can improve gradient flow but does not guarantee
unlimited retention or eliminate every optimization difficulty. The direct
cell-state path contains products of forget gates when their other
dependencies are held fixed; the full derivative is more complicated.

Choose illustrative values $c_{t-1}=0.8$, $f_t=0.9$, $i_t=0.2$,
$g_t=-0.5$, and $o_t=0.7$. Then
\[c_t=0.72-0.10=0.62,\qquad h_t=0.7\tanh(0.62)\approx0.386.\]
These are hand-chosen teaching values, not training outputs. If $f_t=1$,
new content can still change the cell. Changing $o_t$ changes hidden output
without changing this cell update. Predict those effects before using sliders.

\subsection{Implementation, selection, and reproducibility}
With \texttt{batch\_first=True}, input/output layouts are $(B,T,F)$ and
$(B,T,H)$. Ordinary non-projected hidden and cell states retain layout
$(\mathrm{layers}\times\mathrm{directions},B,H)$. PyTorch concatenates
gate chunks in input, forget, candidate, output order. The notebook
reconstructs a learned cell and compares it numerically with \texttt{nn.LSTM}.
Training updates use training data, validation selects a checkpoint, and
test is evaluated afterward. Save raw data identity, seed, versions, window,
split boundaries, normalization, model state, baselines, and measured errors.
A seed alone does not guarantee bitwise agreement across environments.

\section{Hour 2: why attention is useful and how it works}
\subsection{The fixed-context bottleneck}
A final-state encoder compresses a source sequence into one fixed-width
context. Different output steps may need different earlier details.
Attention makes encoder states addressable and forms a context for each
query. It does not guarantee correct retrieval or unlimited context.
Attention can be combined with recurrence: Bahdanau and colleagues used
recurrent sequence-to-sequence modeling in a 2014 preprint presented in
2015. The Transformer architecture came later. Distinguish a retrieval
mechanism from the entire architecture that uses it.

\subsection{Queries, keys, values, and shapes}
Queries encode a retrieval need. Keys encode matching information, while
the corresponding value rows carry content to combine. For $N_q$ queries,
$N_k$ keys, key width $d_k$, and value width $d_v$,
\[Q\in\mathbb R^{N_q\times d_k},\quad
K\in\mathbb R^{N_k\times d_k},\quad V\in\mathbb R^{N_k\times d_v}.\]
Values need not equal keys or have their width. A natural-language question
is not the same object as the query vector produced by a learned projection.

Additive attention can use
\[e_{ij}=v_a^\top\tanh(W_qq_i+W_kk_j+b).\]
Scaled dot-product attention uses
\[S=QK^\top/\sqrt{d_k},\quad
A_{ij}=\frac{\exp(S_{ij})}{\sum_m\exp(S_{im})},\quad C=AV.\]
Scores compare, row-wise softmax normalizes, and value multiplication
retrieves. Rows are query destinations and columns are key sources.
The context has shape $(N_q,d_v)$.

Under independent zero-mean unit-variance coordinate assumptions, the raw
dot-product variance grows with $d_k$. Scaling by $\sqrt{d_k}$ controls
that growth; it is not cosine normalization. Stable softmax subtracts the
row maximum before exponentiation without changing normalized probabilities.
Before later projections and residuals, the weighted context lies within
the convex hull of the available values.

\subsection{An attention calculation with no hidden steps}
Choose teaching vectors
\[Q=\begin{bmatrix}1&0\end{bmatrix},\quad
K=\begin{bmatrix}1&0\\0&1\\1&1\end{bmatrix},\quad
V=\begin{bmatrix}10&0\\0&10\\10&10\end{bmatrix}.\]
Raw dot products are $(1,0,1)$ and scaled scores are approximately
$(0.7071,0,0.7071)$. Softmax gives weights
$(0.4011,0.1978,0.4011)$. The zero raw score does not give zero weight,
because $\exp(0)=1$. Combining corresponding value coordinates gives
\begin{align*}
C_1&=0.4011(10)+0.1978(0)+0.4011(10)\approx8.0222,\\
C_2&=0.4011(0)+0.1978(10)+0.4011(10)\approx5.9889.
\end{align*}
Use full precision in code. Change only the query and predict the new
ranking. An optional positive temperature divides the scores: below one
sharpens weights, above one softens them. The standard example uses one.

\subsection{Self-attention, cross-attention, and masks}
Self-attention projects one sequence $X$ into $Q=XW_Q$, $K=XW_K$, and
$V=XW_V$. Cross-attention derives queries from one stream and keys/values
from another. An encoder-decoder model uses decoder queries with encoder
memory. A decoder-only model normally has no separate source encoder.

In causal next-token training, input position $i$ predicts token $i+1$,
so it may access itself and earlier inputs: $j\le i$. Set future scores
to $-\infty$ before softmax, yielding zero weights. Padding masks hide
nonexistent tokens instead. Combine permissions and exclude padded query
outputs from loss. An all-masked row has no valid normalization; prevent
it or specify a deliberate special case. Mask conventions differ across
APIs. PyTorch scaled-dot-product attention uses Boolean True for allowed
entries, whereas some other attention calls use True for blocked entries.

\subsection{Benefits, costs, and interpretation}
One dense attention layer provides a direct query-to-key computational
path to a distant position. Known query operations can run in parallel.
Dense score interactions involve $T^2$ pairs per head: doubling context
length quadruples pair count. Exact fused implementations can reduce
intermediate memory without changing dense pair-count scaling.

Attention weights describe routing within the computation; they do not
prove causal importance. Keep $Q$ and $K$ fixed and change $V$: weights
stay identical but the result can change. Removing and renormalizing a
source is a different intervention. Explain what a particular experiment
supports rather than interpreting an attractive heatmap as a complete
causal explanation.

\section{Hour 3: assemble and test a Transformer}
\subsection{Tokens and position}
Token IDs index learned vectors; their numeric magnitude has no inherent
semantic distance. The teaching model has eight symbolic tokens, width 32,
four heads, and context length 16. Vocabulary size and embedding width differ.
Unmasked attention without positional information is permutation-equivariant:
permuting input rows permutes outputs correspondingly. Order must therefore
be represented deliberately. Our model adds learned absolute position
vectors. The original Transformer used sinusoidal encodings; other position
mechanisms also exist. A causal mask already introduces order constraints,
so do not apply the unmasked equivariance claim unchanged to a masked model.

\subsection{Heads, nonlinear transformations, and residual paths}
Each head uses learned projections and its own routing weights. Concatenate
head outputs and apply a learned output projection. At fixed total width
$D$, $H$ heads have width $D/H$ in this implementation. More heads do not
automatically multiply the total width, and heads have no assigned semantic
roles. Attention mixes positions, while a shared tokenwise MLP transforms
features independently at each position. The notebook expands to $4D$,
applies GELU, and projects back to $D$.

Residual addition requires matching dimensions and preserves an identity
route. The teaching decoder uses pre-normalization:
\begin{align*}
u&=x+\operatorname{Attention}(\operatorname{LN}(x)),\\
y&=u+\operatorname{FFN}(\operatorname{LN}(u)).
\end{align*}
The 2017 paper used a full encoder-decoder with post-normalization and
ReLU. Our one-block decoder is an explicitly identified educational
variant, not a verbatim copy of that architecture. Its causal self-attention
has no encoder-decoder cross-attention sublayer.

\subsection{Objective, baseline, and generation}
Shift targets one token ahead. Logits have shape $(B,T,V)$, and
cross-entropy minimizes
\[\mathcal L=-\frac1{BT}\sum_{b,t}\log p_\theta(y_{b,t}\mid x_{b,\le t}).\]
Pass logits directly to \texttt{nn.CrossEntropyLoss}. Training computes
many known positions in parallel under the causal mask. During generation,
the model predicts one unknown token, appends it, and repeats. Parallel
training does not imply arbitrary autoregressive generation is parallel.

Separate synthetic streams use seeds 2026, 2027, and 2028 but share one
motif with small random perturbations. This is generalization to new
realizations of that process, not transfer to a different engineering
domain. A training-only smoothed bigram model is the reference. Uniform
prediction over eight tokens has NLL $\log 8\approx2.079$ nats/token.
Select a checkpoint by validation NLL; evaluate test afterward.

\subsection{Behavioral checks and limits}
Plot an actual trained attention head with query rows and key columns.
Upper-triangle zeros are mask constraints, not learned discoveries.
Change tokens only at positions 8 onward and compare logits at positions
0--7. They should remain unchanged within the declared tolerance.
This tests one concrete causal behavior, not every deployment property.
Separately test permutation equivariance of the unmasked operation.
Removing positions from an already trained decoder is an inference
intervention, not a fair separately trained architecture comparison.

The model generates symbolic states, not evidence-backed natural language.
A token generator is only one possible component of an agent. An agent
also needs a bounded objective, state handling, validated tools,
permissions, stopping rules, and evidence review. Generation alone does
not provide factual verification or controlled execution.

\section{Completion and explain-back}
H1 must produce a baseline table, residual plot, a reconstructed learned
cell update, and evidence JSON. H2 must produce scores, weights, context,
mask assertions, and an interpretation-limiting intervention. H3 must
produce a trained-model comparison, head plot, future-token check, and
configuration record. Predict the effect of one change before executing,
keep the original evidence, and finish with a claim, evidence, limitation,
and a test that could challenge the claim.

\section{Primary sources}
\begin{itemize}
\item Hochreiter and Schmidhuber (1997), \emph{Long Short-Term Memory}:
\url{https://doi.org/10.1162/neco.1997.9.8.1735}.
\item Bahdanau, Cho, and Bengio (2014/2015), recurrent attention:
\url{https://arxiv.org/abs/1409.0473}.
\item Vaswani et al. (2017), \emph{Attention Is All You Need}:
\url{https://arxiv.org/abs/1706.03762}.
\item Dive into Deep Learning:
\url{https://d2l.ai/chapter_recurrent-modern/lstm.html} and
\url{https://d2l.ai/chapter_attention-mechanisms-and-transformers/transformer.html}.
\item Executable course notebooks:
\url{https://github.com/CHjonnyYu/MLES-Fall-2026}.
\end{itemize}

\end{document}
